\section{Polymorphism}
Difference in code behaviour without changing existing code. Allows succinct,
future-proof code.
\paragraph{Dynamic Binding}
Determining which method to be invoked from RTT of object.
\verb|x.invokedMethod(params...)|
\begin{enumerate}
  \item During compile time
    \begin{enumerate}
      \item Determine method descriptor of invoked method with \verb|CTT(x)|
      \item Search \emph{ALL} methods in \verb|CTT(x)| that can be invoked
    \item Choose the \emph{most specific} method.\ \textcolor{Bittersweet}{If none exists, error.}
      \item Store method descriptor in the bytecode
    \end{enumerate}
  \item During run-time
    \begin{enumerate}
      \item Retrieve method descriptor from bytecode
      \item Search for first matching method:\\
        $\Rightarrow$ Starting from \verb|RTT(x)|, iterating to parent class
    \end{enumerate}
\end{enumerate}
\textcolor{Blue}{Most Specific Method}: If arguments to \verb|M| can be passed
into \verb|N| without compilation error, \verb|M| is more specific than
\verb|N|

\textcolor{Bittersweet}{Class Methods}: No Dynamic Binding. Step 1 taken and
CTT's method is executed.
